\section*{Erd\H{o}s Problem 696: prime chains and divisor chains}

\paragraph{Statement and claim.}
Let \(h(n)\) be the longest increasing chain of prime divisors
\(p_1,\ldots,p_h\) of \(n\) satisfying \(p_i\mid p_{i+1}-1\).
Define \(H(n)\) in the same way with arbitrary positive divisors.
We prove, for almost all \(n\),
\[
 h(n)\asymp H(n)\asymp\log_* n,\qquad H(n)/h(n)\leq18.
                                                               \tag{1}
\]
In particular the proposed almost-everywhere divergence of the
ratio is false. This is a reconstruction of the coarse resolution
reported at \url{https://www.erdosproblems.com/696} and its forum,
using the prime-successor method in the manuscript hosted at
\url{https://github.com/davidturturean/erdos-696}.
That manuscript attributes its stronger estimates to work with
GPT-5.5. No new resolution or sharp normal-order constant is
claimed here.

\paragraph{External analytic inputs.}
We use the classical Siegel--Walfisz theorem in the following
fixed-exponent form: for coprime \(a,q\), uniformly for
\(q\leq\log t\),
\[
 \pi(t;q,a)=\frac{\operatorname{Li}(t)}{\varphi(q)}
             +O(t e^{-c\sqrt{\log t}}).                         \tag{2}
\]
See A. Walfisz, \emph{Zur additiven Zahlentheorie. II},
Math. Z. 40 (1936), 592--607; the cited manuscript states
and applies the same uniform input.
We also use the elementary classical bounds
\(\sum_{p\leq z}\log p=O(z)\) and
\(\sum_p1/p=\infty\).
All other probability, counting, and tower estimates are given below.

\paragraph{Few short successor steps in a typical divisor chain.}
Call a pair bad if \(d\geq Y\), \(d<e<\exp(\sqrt d)\), and
\(e\equiv1\pmod d\). Then \(\gcd(d,e)=1\). For fixed \(d\),
writing \(e=1+jd\) gives
\[
 \sum_{\substack{d<e<\exp(\sqrt d)\\e\equiv1\ (d)}}\frac1{de}
 \leq\frac1{d^2}
       \sum_{1\leq j\leq \exp(\sqrt d)/d}\frac1j
 \ll d^{-3/2}.
\]
The number of \(n\leq x\) divisible by both members of any bad
pair is therefore, by a union bound,
\[
 \leq x\sum_{d\geq Y}\sum_e\frac1{de}\ll xY^{-1/2}.             \tag{3}
\]
For every other \(n\), consecutive chain terms above \(Y\)
satisfy \(d_{i+1}\geq f(d_i)\), where \(f(t)=\exp(\sqrt t)\).
For all sufficiently large \(t\),
\[
 f(f(t))=\exp\bigl(\exp(\sqrt t/2)\bigr)\geq\exp(t).
\]
Define \(T_0=e\), \(T_{j+1}=\exp(T_j)\), and let
\(L=\min\{j\geq0:x\leq T_j\}\).
This differs from any usual convention for \(\log_*x\) by
at most a fixed constant. Every two chain steps above \(Y\)
raise the tower level by at least one. There are at most
\(2L+O(1)\) such terms, and at most \(Y\) terms below \(Y\).
Taking \(Y=\lfloor\sqrt L\rfloor\), we obtain
\[
 H(n)\leq3L
 \quad\hbox{for all but }o(x)\hbox{ integers }n\leq x.            \tag{4}
\]

\paragraph{Enough prime successors in a very large window.}
For large \(y\), put
\[
 Z(y)=\exp(\exp y),\qquad F(y)=\exp(\exp(y^2)).
\]
For every prime \(p\leq y\), partial summation in (2) gives,
uniformly in \(p\),
\[
 \sum_{\substack{Z(y)<q\leq F(y)\\q\ {\rm prime}\\q\equiv1\ (p)}}
       \frac1q
 =\frac{y^2-y}{p-1}+O\bigl(e^{-c'\sqrt{\log Z(y)}}\bigr)
 \geq y/3.                                                     \tag{5}
\]
Indeed, the modulus condition holds throughout the interval.
The integrated error is bounded by a constant times
\(\int_{\sqrt{\log Z(y)}}^\infty u e^{-cu}\,du\), and the
main integral is
\((\log\log F(y)-\log\log Z(y))/(p-1)\).
Thus no constant in (5) depends on the chosen prime \(p\).

\paragraph{Independent prime choices and transfer to integers.}
Let \(y_0=L\), \(y_{j+1}=F(y_j)\), and
\(K=\lfloor L/5\rfloor\). Consider independent Bernoulli
variables indexed by primes \(q\leq y_K\), with success
probability \(1/q\).
With probability tending to one, there is a successful
prime \(p_0\leq L\), since the probability of none is
\(\prod_{p\leq L}(1-1/p)\longrightarrow0\).
Having chosen \(p_j\leq y_j\), seek a successful prime
\[
 Z(y_j)<p_{j+1}\leq y_{j+1},\qquad p_{j+1}\equiv1\pmod{p_j}.
\]
The window is disjoint from every window previously inspected.
Conditional on all previous choices, (5) shows that the
probability of failure is at most
\[
 \prod_q(1-1/q)\leq\exp\left(-\sum_q1/q\right)
 \leq e^{-y_j/3}\leq e^{-L/3}.
\]
A union bound over \(K\) stages now gives a prime chain of
length \(K+1\geq L/6\), with probability \(1-o(1)\).

The size of the prime cutoff matters for transferring this
model to uniform integers \(n\leq x\).
Let \(r_0=\min\{r:L\leq T_r\}=o(L)\).
For large \(y\), \(y^2\leq e^y\), so \(F(y)\leq\exp(\exp(\exp y))\).
Inductively,
\[
 y_K\leq T_{r_0+3K}\leq T_{L-3}<\log\log x
\]
for all sufficiently large \(x\); the last inequality uses
\(x>T_{L-1}\).
Set \(Q=\prod_{p\leq y_K}p\). The classical bound quoted above
gives
\[
 Q\leq\exp(O(y_K))\leq(\log x)^{O(1)}=o(x).
\]
On a uniformly random residue modulo \(Q\), the divisibility
events for its primes are exactly the independent model,
by the Chinese remainder theorem.
For every event depending on that residue, its probability
on \(\{1,\ldots,\lfloor x\rfloor\}\) differs by \(O(Q/x)=o(1)\).
Consequently
\[
 h(n)\geq L/6
 \quad\hbox{for all but }o(x)\hbox{ integers }n\leq x.            \tag{6}
\]

\paragraph{Conclusion and scope.}
Combining (4), (6), and \(h(n)\leq H(n)\) gives
\(H(n)/h(n)\leq18\) outside \(o(x)\) exceptions up to \(x\).
After discarding the \(O(\sqrt x)\) integers below \(\sqrt x\),
the tower ranks of \(n\) and \(x\) differ by at most a constant.
Thus (4)--(6) also prove the two order-of-growth assertions in (1).
This settles the stated coarse estimation and ratio question,
relative to (2). It does not reproduce the sharper constants
discussed in the cited manuscript.

\paragraph{Completion estimate.}
\textbf{COMPLETION: 100\%}
